\subsection{人类作者的惯用词、段落骨架与行文节奏}

从 A217、A001、A022、A171、A0127、A0165、A0175、A092、A016、A053、A163、A178、
A242、A196、B007、B026、B050、B160、B030、B035、B086、B226、B311、B477、B159、
B195、B196、B060、B157、C006、C085、C169、C283、C155、C229、C050、C126、C228、
C038、C063、C094、C234、C023 与 C132 共四十四篇逐页细粒度文风样本，并结合其余
十五篇结构化动作卡看，“有人味”不等于
口语化，也不等于把连接词替换一遍。
它来自作者在段落里留下真实的观察和取舍。A217 常用“注意到……，故……”把几何
特征接到降维推导，用“考虑到……，本文将……”说明为何采用某种坐标或计算口径；
B007 常以“观察附件……发现……，利用……判断……”交代数据处理起点，以“由上述
分析可知……，故……”从比较结果转入模型改进。这样的“故”前面必须有可检查的
事实，后面必须落到一个实际动作，不能只连接两句结论。
B026 还常把“但由于……所以我们预测……因此，我们建立……”用于散点和机理上界
之间的推理，并用“但经数据比对，我们发现题目所给出的数据中没有……”引出补测；
这类长承接只在前后事实确实相反或存在缺格时使用。
B050 的常用动作不是增加连接词，而是把限制条件直接嵌入裁决句：例如先承认四次
多项式的 $R^2$ 达到 1，紧接着写明样本只有五个点，因而不把插值结果等同于整体拟合；
又在切换方法前保留旧方案“结果不理想”的具体对象。写作时可复用“虽然……，但由于
……，故选取……”“在原模型的基础上加入……，调整后 $R^2$ 由……变为……”和“考虑
到上述规划结果……，下面改用……”等推进关系，但省略处必须填本题读数、限制和动作。
B160 进一步给出更细的比较和诊断承接：“当……相同、……相同，只改变……时，将……
组进行对照”先证明组间可比；“由图……可以看出……；在同一温度下……；从整体来看……”
把两个响应分开读取；“从回归的相对效果看……；从绝对效果看……；从方差分析方面看……”
则按指标功能安排句序。这些表达之所以自然，是因为每个分句都填有保持量、组名、读数或
统计量，而不是因为用了更多连接词。

本轮补入的十五篇又把“人味”落到了模型接口上。A070、A147、A195、A212 的共同点是
先把机理候选、控制体内核、分层目标或基线残差写清，再说明哪一项证据触发降阶、增量、
第二层裁决或局部修正；A028、A115 则把近共线量级和硬约束超限直接写成几何简化或目标
改写的理由。B078、B108、B125、B175 不从“综合决策”起笔，而是先区分信息集、在线/离线
状态、库存接口和玩家目标，只有信息边界改变后才引入随机转移或 MDP。C109、C142、C170、
C227、C305 的段落都保留字段缺口和模块交接：评级或违约标签缺失时先补输入，风险排序
再进入额度、利率和放贷决策，疫情只改影响层。可调用的短句是“问题二的目标没有变化，
缺口只在……；先补……，再把输出接回……”，以及“前问离线路径不能直接复用，因为当前
状态只保留……”。这些句式的前提是本题确实存在相同接口和已点名的缺口，不能用来装饰
没有跨问依赖的段落。

B030 的惯用承接更适合几何分类和长仿真结果。“根据……将相对位置关系分为如下三类”
先保证覆盖；“但在求解过程中，发现……在几何上具有对称性，可通过旋转重叠”再给压缩
依据；“固定……，遍历……，记录所得全部结果，并绘图输出”把程序动作改写为可核的论文
步骤。需要复用时，作者先写“由问题一（1）中的定位模型4可知”或“联合式……得”，
而不是直接说“同理可得”。这些词句只在分句中填入具体形态、变量、范围、式号和输出时
才有人味；留下省略号的空壳仍是模板腔。

B035 又补出一组更接近日常推导的转折词。作者用“事实上”“不难发现”“注意到”“同理”
承接几何事实，但每次后面都紧跟圆、弦、交点或角度区间；需要换方法时，先说原来四区域
分类为何麻烦，再写“这里我们应用一个相对更简单的求解方法”。误差论证用“仍旧以……
为例”“下面我们通过……进行说明”把例子、图形和区间表逐步接起。结果比较不回避局部
反例，而以“在局部上可能……占优势，但从全局上看……”转入误差较大点先被修正的机制。
这些惯用词可用，但只有补入当前题目的对象、数值和下一步动作时才不同于 AI 式空转。

B086 的惯用承接更适合处理信息边界、算法比较和几何流程。作者先用“该问本质上是一个
……问题”给小节定性，再以“考虑到……故需……”把限制接到动作；局部搜索处明确写
“由于……信息不共享，故无法通过全局……”，随后用“注意，此处定义的目标函数不是
全局函数”阻止概念混用。定义事后评价量时又补一句“这是合理的，因为……并没有根据
全局误差函数做出调整”，把评价和决策分开。预处理证明后使用“由此可见……更重要的
是……”从误差界转到搜索次数；方案裁决前使用“特别注意……无实际参考价值，其仅用于
比较……”限定资源代理。稳健性段以“需要注意的是，单凭一组原始数据并不能准确检验
……故……”引出十组扰动。这些表达的人味来自它们公开了限制、用途和下一步，而不是
来自句首词本身。

B226 的惯用承接集中在“先归约、再命名、后推导”。“由问题一知，只要已知……和……，
就可以求解……，因此我们需要……”把复用接口写得很清楚；“为了让模型的建立过程更加
直观易懂，我们记……”让面、线和向量先拥有共同称谓；“根据……规程……出于建模的
严谨性，我们在下文证明……”把经验规范和数学授权分开。贪心求解用“以……为始边”“覆盖
宽度恰好能达到……”“以 1 m 为步长”逐项写出程序动作；求解器切换则先说三个未知参数
及其较大范围会使循环遍历耗时，再引出粒子群。分区段主动使用“我们人为地将……”并列出
依据，拟合不佳时又以均方误差触发局部再分。结果段采用“区域 V 的测线极多。这是因为……
又因为……这就需要……”的区域特定因果链。这些句式只有填入旧模型、两个缺失量、几何
对象、边界、步长、误差指标和异常区域时才可复用；若只留下“因此”“这是因为”，仍是 AI
式模板。

B311 的自然感首先来自它不掩去定义上的犹豫。作者先承认“由题意可以得到两种理解”，
分别推到可计算结果，再用共同数值例和 CAD 图中的可观察差异作裁决；“模型一”“模型二”
不是并列堆砌，而是一真一假的候选。跨问处的“比较本问与问题一可知”后面则紧跟不变
结构、改变的水深和夹角，以及需要替换的函数输入。小坡度也没有被一句“影响较小”略去，
而是先换算到 2 海里，看到累计差值后再决定方向。方向选定以后才补一句与测量规范一致，
把题内判断和外部印证分开。

本篇还说明，有序连接词并非天然带有 AI 腔。“首先、其次、再次、最后”在原段中分别
检查对称、单调、特殊方向不变和极限式，四个分句的主语与裁决对象都不同，因此读者
能逐项核对。递推段用“从西边界开始”“当第 34 条……”“故第 35 条舍弃”完成对边闭合；
最不利位置段则先找等深线间距最大处，再写“若该处完全覆盖，其余位置亦可覆盖”。可
迁移的不是这些句首词，而是“条件--数值--边界--动作”均不省略的组织方式。

B477 的惯用承接常出现在作者发现原路走不通的时候。“……的结果是反直觉的，若继续
使用，……将失去意义”先把旧定义的失效对象说透，随后才写“为适应……，将……改为
……”。“由于前述流程尚未确定从浅处还是深处起算，两种情况均作计算”则公开未定选择，
再以同一越界面积裁决。方法切换处的“这种做法在后续计算中有以下几点负担”后面列的是
四个可核对象，而不是“效率较低、效果不佳”的同义反复；新方法紧接着给出深度、梯度、
坡度和覆宽四个下游量。路线段的“经过初步尝试，发现……”也保留起点依赖与漏测，继而
用“因此只对……”限定返工范围。可迁移的骨架是“具体现象--受影响对象--必要修改--
交付接口”，不能只摘取“反直觉”“针对上述问题”“经过尝试”这些句首词。

C006 的段落推进更接近真实业务建模。作者先承认“供货连续性”不能直接读取，再找间隔和
连续周数作代理；在周能力估计处不用历史均值，不是泛称“考虑周期性”，而是分别写出均值
过于均匀和作物具有周期两个理由。遇到“优化结果没有材料偏好”这一反常现象时，段落也没有
立即换算法，而是先检查成本尺度和等产能关系，再局部修改 A/C 约束。结果部分把 24 周总体
走势与两个具体周次并列，并用“最大约降 $5\%$、平均仅降 $0.956\%$”限定改进强度。这些句子
自然，是因为它们允许“不采用”“不明显”和“继续沿用”出现，并各自给出理由。

C085 的人类痕迹集中在“具体失败--局部修补”和“先改比较口径”两种动作。作者不写
“为提高算法性能引入动态权重”，而是先说全局固定权重会造成“总体情况良好但是某一周
极度缺少供货”，再按该周到期望产能的差更新下一状态权重；也不泛称“考虑周期性”，
而是点出第 7、8 周连续淡季，由此取两周前视。类似地，“历史成本略低”没有直接成为
优劣结论，因为两套方案的产能不同；换成成本/产能后才作裁决。这里可复用的惯用承接是
“本文设定了两个规则……在以上两个规则都满足的情况下……”“若全局只使用……可能会
导致……因此……”“经过数据分析，发现第……周……因此……”和“……虽然略低，但由于
……未达到……故取……作为判断标准”。每个省略处都必须填阈值、周次、基准和数值。

C169 的惯用词服务于“先给理由，再决定纳入或舍弃”。构造非对称指标时，作者连用“为了
体现出……的影响明显重于……”“由于……的正负造成的影响不同，所以……满足下列条件”
和“综合考虑以上条件和生活实际，构造出……”，顺序是业务差异、数学条件、具体函数。
发现供货规律后又写“容易发现……但在下文的分析中不单独分析，原因有如下两点”，允许
分析以舍弃结束。初始状态用“考虑到企业并非刚刚成立……因此不能简单地认为……为 0”，
可分性用“由于……无记忆……所以……；为了简便计算可以一周为单位考虑”。这些表达像人写，
不在于连接词本身，而在于每个连接词后都跟着可核对的业务事实和建模后果。

C283 的自然感更多来自作者把犹豫、否决和边界留在正文中。假设段不是直接宣布简化成立，
而是写“由题可知……；若假设……，则……难以满足……，故原假设不成立”；权重段用“但根据
实际操作结果来看……与事实不吻合，需要在适度范围内修正”承接客观结果与业务判断。算法段
先分别写 LP 和 DP 的具体代价，再落到双层贪心；达到库存边界后用“见好就收”收束，并紧接
继续增加订单的损失。随机结果用“多次实验”“频率较高且目标较优的典型代表”限定结论强度；
方向权重用“一方面……不可过小；另一方面……不可过大，否则……”同时写出下界和上界。
这些惯用词不能脱离对象照抄，真正应复用的是“提出判断--暴露反例或代价--限定范围--作出选择”。

C155 的自然感来自作者愿意写出“为什么这里不做”。对颜色缺失，正文不是机械插值，而是先说明
相同已知条件仍出现多种颜色，再指出随机表面取样不能代表整件文物，随后落到“因此本文认为不能
通过已知数据填补……，后文涉及颜色时忽略……”；对表单 3，又用“与表单 2 不同的是……”说明
风化字段已经可直接观察，因而删除重复估计步骤。卡方段常写“由表……可知……满足使用 Pearson
卡方检验前提，因此……；……不满足……，因此……”，把检验名称放在前提之后。回归段以“在后续
预测过程中，仅对 $R^2$ 较高的成分进行预测，对较低成分采取不变策略”把指标直接接到动作；异常
距离段则按“当回归方程确定时，距离 $d$ 是唯一决定变量--异常值会导致预测失真--为了减小影响而
压缩距离”推进。可复用的不是句首“因此”，而是每个因此之前都有字段、样品编号、阈值或失真机制，
之后都有剔除、保留、换检验、预测或不预测的具体动作。

C229 的自然感则来自作者把同一字段内部的证据差异写出来。颜色填补不是先定方法再套数据，而是
先写“按此分为两个类别分别进行填补”，逐一列出 40 对 39、58 对 11 的参照关系；19、48 号
因候选化学样品“参考性较低”，又改用众数。近似零值段用“不能单纯将缺失值插补为 0，可看作
近似零值”承接检测限和舍入语义；作图段直接承认“因二氧化硅占比较大，图中不能很好地体现其余
成分”，故保留全图并补画去除图。亚类段先写“没有已知因变量，因此该类问题为无监督学习问题”，
结果不清时又保留“仅可粗略将其分为三类”。可复用的是“候选质量触发方法分流--图形局限触发
补图--有无标签决定方法身份--证据强弱决定语气”，不是热卡、聚类或“故”字本身。

C050 的惯用承接更接近真实的数据分析过程。正文先写“考虑到数据较为庞大和复杂，为更好地分析
和理解数据”，随后不是泛泛宣布预处理，而是逐项观察丰富度、销量、损耗和分布；对象过多时又写
“由于……种类较多，本文选取……”并立即交代代表规则和附件去向。回归段的“由回归系数可知”
后面紧接 $1\,\mathrm{kg}$ 与价格变化的业务读法，特征段的“为有效判断……，本文选取……”后面
则给出对象、口径和代理含义。这些句子有人味，是因为作者在每次缩减数据、选择对象或解释系数时
都公开了判断依据，而不是因为“考虑到”“由于”“由此可知”出现得多。

更值得保留的是作者没有把代理量写成唯一解释。打折次数既可能对应受欢迎和供不应求，也可能对应
销售较弱或保鲜期较短；文章先并列竞争机制，再要求补看销量、反馈与竞品信息，最后分别落到补货、
定价或促销动作。现有附件与待采外部数据也分别成段：前者支持已经量化的关联，后者只说明采集方式、
可辨机制和决策用途。可迁移的段落不是“还需采集更多数据”这一空句，而是“现有证据能回答什么--
哪两种机制仍无法区分--补哪项观测--不同结果分别采取什么动作”的完整判断轨迹。

C126 的惯用承接更像作者在一边处理数据、一边向读者说明为什么这样处理。正文先用“由于……，
直接采用……会……”保留不正态、零值和计算负担，再写“为避开……同时降低……，用同期平均法
将……”；“为此确定”后面跟的是新时间序列和迭代设置，不是重复结论。对象太多时，作者常写
“这里仍以……为例进行说明”，并在下一句交代大矩阵为何不宜展开；参数表之前又用“因篇幅限制，
这里只给出……”明确正文保留什么、支撑材料保存什么。这些句式自然，是因为每个省略处都要填
数据现象、计算代价、代表对象和实际文件，而不是把“由于”“为此”“因篇幅限制”当成润色词。

结果段的承接更值得模仿。“由实际经营过程可知”后面紧跟隔日难售这一事实，并据此确定 VAR(2)；
“表中第一行……第二行……第三行……”先教读者读取系数、标准差与检验量，才综合 AIC、SC、
$R^2$ 和 F 值把结论限定到未来一周；“注：负数表示……”则把负预测转成余货、折价和不补货。
最后，“从另外一个方面进行筛选”不是平行堆模型，而是用半小时销售频率独立核对 29 个对象。
这组表达真正可复用的，是“事实--尺度--选型”“表项--指标--限定结论”“数值--状态--动作”和
“主路线--独立规则--有限一致”四条判断轨迹。

C228 的惯用词同样依附于具体判断。开放题先写“由于题目未明确给出具体的探究方向，需要先考虑
从哪几个角度入手”，随后立即列出时间、销售和关系三个观察面；这里的“由于”不是装饰，而是在
解释为何正文先聚焦而不是直接建模。预测段用“首先采用……，结果发现效果较差，因此……”保留
候选方法的失败状态；单品段则写“观察最近筛选出的……，发现……”后逐项指出缺失、无明显趋势
和随机波动，再把复杂模型降为加权平均。两段都让“结果发现”后面跟可观察事实，让“因此”后面
跟可执行选择。

求解器切换时，作者先写“与问题二不同”，再列出本问新增的 $0/1$ 进货标志、连续决策和双目标，
最后才交代混合编码与 NSGA-II；效果解释则先说明为什么选早七月的历史经营区间，再把方案利润放入
这个基线中。改进建议用“尽管……直观便捷，但其主要从……出发”承认成本导向定价的作用，再从
其边界推出竞品、预期价格和收入数据。应迁移的是“开放任务--观察面”“候选效果--换路”“数据
不足--降级”“变量结构--求解器”“参照理由--有限比较”和“现有方法--边界--补采”的段落次序，
不是机械复制“由于、首先、因此、尽管”。

A0127 的自然感来自作者愿意在算法之前交代工程判断。开头用“它实际上是一道优化问题”给全题定性，用“问题二是问题一的反问题”说明问间关系；随后不是笼统说计算困难，而是先把题面公式中可直接得到的量剔除，只留下五类效率。塔影段的“考虑到”后面有明确负担、数量关系和中心对称，因而“以镜面中心作判断”是一个有来源的近似；布场段的“从图中可以看出”后面紧接北侧高效、镜面北移和塔位南移三个动作；问题三的“若直接采用遍历法……存在计算量过大的问题”又自然引出先二分、再变步长。灵敏度段还保留“快速增加至一定值后几乎不变”的饱和描述。这里应模仿的是“事实—负担—取舍—动作—边界”的信息顺序，而不是复制几个句首。

A0165 的惯用表达更像作者在顺着计算过程说话。摘要常先写“将……划分为……，分别进行建模”，随后用“通过……确定……”“根据……计算……”“利用……判断……”逐级交付中间量；问题分析则从“为计算……首先需要……”起笔，先补齐太阳位置和镜面法向，再进入遮挡、截断和功率。作者写简化时不只说“为提高效率”，而是交代“分析其中一个变量时，将其他变量固定”，并用完整计算作基准说明 $2\times2$ 栅格的偏差。图形段的“由第一问散点图可知”后面紧跟塔向南移动这一设计动作；问题三的“在问题二的基础上”后面则列出不再改变的宽度、塔位和位置，只把重心放在各圈高度。检验段先说明前两问由第一问变化而来，再改变采样时刻核对峰值移动是否符合太阳视运动。应模仿的是每个常用词后都有具体对象、冻结量、动作和裁决，不把“通过、根据、因此”留下空壳。

A016 常用“由于……因此应当满足……”“同时由于相邻……只需要……，因此选取……”“若……则撤回上一步，并将步长除以 2”“但是由于我们按照时间遍历……立即停止，因此……”推进。这里连续出现的“由于、因此、但是”并不显得模板化，因为每个词都在承担不同责任：第一个引入物理顺序，第二个裁决多根，第三个规定数值回退，第四个用事件先后排除后发分支。“可以发现……唯一确定，因此……固定不变”则用来从几何关系推出没有可优化自由度；“不妨设……变为 $k$ 倍，由……可知……也变为 $k$ 倍”把局部比例传到全链。真正该写进 Skill 的不是这些词的表面频次，而是连接词后必须出现条件、被选对象、执行动作和结论边界。

